% !TeX root = ../main.tex
\section{Extended Lagrange--d'Alembert principle in virtual-power form}
\label{sec:virtual_power_derivation}

The governing equations are derived here from an extended
Lagrange--d'Alembert principle written in virtual-power form.  Its time-integrated
constrained structure follows the extended Hamilton principles developed by
Bedford and Drumheller for porous media and reacting immiscible mixtures, in
which constituent motions, volume fractions, and conversion variables are
varied together
~\cite{bedford1979variational,drumheller1980thermomechanical,bedford1983theories}.
The organization of its integrand into internal, external, inertial, and
constraint virtual powers follows the continuum-mechanical method of virtual
power, whose classical statement is instantaneous
~\cite{germain1973virtualpower,gurtin2010continua}.  The time integral converts
these virtual powers into virtual work and permits temporal integration by parts
for accumulated state coordinates.  Reaction progress and relative component
fluxes are included as independent dissipative process rates; their constitutive
force--rate relations are subsequently restricted by an Onsager dissipation
potential~\cite{doi2011onsager,wang2021onsager}.

The mixture contains arbitrary sets of solid phases \(\mc S\) and fluid phases
\(\mc F\).  The index \(\xi\in\mc S\cup\mc F\) denotes a phase, \(\alpha\)
denotes a component in that phase, and \((m)\) denotes a conversion mechanism.
Each phase has its own motion \(\mbf x_\xi\), velocity \(\mbf v_\xi\),
deformation gradient \(\mbf F_\xi\), and Jacobian
\(J_\xi=\det\mbf F_\xi\).  Its current component density is
\(\rho_\xi^\alpha=\phi_\xi\bar\rho_\xi\eta_\xi^\alpha\).
The component storage multiplier is \(\pi_\xi^\alpha\), the phase
composition-normalization multiplier is \(\Pi_\xi\), and the component-rate
constraint multiplier is \(\mu_\xi^\alpha\).  The entropy inequality in
\sref{MC_residual_dissipation} identifies this same multiplier as the
thermodynamic chemical potential.  All quantities are considered over an
arbitrary interval \(t\in[t_1,t_2]\).  The accumulated conversion and transport
masses, \(\mc C_\xi^\alpha\) and \(\mc D_\xi^\alpha\), are primitive state
coordinates.  Their rate variations are therefore compatible variations,
\(\delta\dot{\mc C}_\xi^\alpha
=D_\xi(\delta\mc C_\xi^\alpha)/Dt\) and
\(\delta\dot{\mc D}_\xi^\alpha
=D_\xi(\delta\mc D_\xi^\alpha)/Dt\), rather than independent virtual rates.
The accumulated-mass variations are taken to vanish at \(t_1\) and \(t_2\),
which later permits temporal integration by parts without endpoint terms.
When electrical effects are retained, the electric potential \(\varphi\) is also
varied.  In the electroquasistatic specialization the electric field is
\(\mbf E=-\nabla_{\mbf x}\varphi\).  The corresponding virtual-power term is
written with the virtual potential rate \(\delta\dot\varphi\).  A polarization
field \(\mbf P\) may either be determined by reversible stationarity or treated
as an internal variable with its own dissipative evolution.

Let \(\delta\mbf v_\xi\) be an arbitrary virtual phase velocity.  The symbols
\(\delta\dot\phi_\xi\), \(\delta\dot{\bar\rho}_\xi\),
\(\delta\dot\eta_\xi^\alpha\), and, when electrical effects are retained,
\(\delta\dot\varphi\) denote independent virtual state rates.  In
contrast, \(\delta\mc C_\xi^\alpha\) and \(\delta\mc D_\xi^\alpha\) vary the
accumulated state coordinates described above.  The independent dissipative
process-rate variations are
\(\delta\dot r_{(m)}\), \(\delta\mbf j_{\xi,{\rm disp}}^\alpha\), and
\(\delta\mbf j_{\xi,{\rm diff}}^\alpha\).  Multiplier variations enforce the
four constraints.  The extended Lagrange--d'Alembert principle in
virtual-power form is
%
\begin{equation}
	\int_{t_1}^{t_2}\left(
	\delta\mc P_{\rm in}
	+
	\delta\mc P_{\rm ext}
	+
	\delta\mc P_{\rm rev}
	+
	\delta\mc P_{\rm diss}
	+
	\delta\mc P_{\rm elec}
	+
	\sum_{k=1}^4\delta\mc P_{C_k}
	\right)\dt
	=0,
	\label{eq:vp_principle}
\end{equation}
%
for every admissible collection of these virtual fields.  The subscripts
\({\rm in}\), \({\rm ext}\), \({\rm rev}\), \({\rm diss}\), and \({\rm elec}\) denote inertial,
external, reversible, dissipative, and electrical virtual power, respectively, while
\(C_k\) denotes the four constraint contributions.  The statement is
time-integrated as in the Bedford--Drumheller constrained Hamilton formulation,
whereas each term in its integrand is expressed as a virtual power.  The
electrical contribution is absent in the neutral specialization.

\subsection{Reference mass measure and virtual rates}

For component \(\alpha\) of phase \(\xi\), the augmented mass per unit phase
reference volume is
%
\begin{equation}
	\rho_{\xi0}^\alpha
	+
	\mc C_\xi^\alpha
	+
	\mc D_\xi^\alpha
	=
	J_\xi\rho_\xi^\alpha.
	\label{eq:vp_reference_component_measure}
\end{equation}
%
Here \(\rho_{\xi0}^\alpha\) is fixed, \(\mc C_\xi^\alpha\) is accumulated mass
added by conversion, and \(\mc D_\xi^\alpha\) is accumulated mass added by
relative component transport.  Consequently, the virtual rate of mass addition
to this referential measure is
\(\delta\dot{\mc C}_\xi^\alpha
+\delta\dot{\mc D}_\xi^\alpha\).  Dividing this sum by \(J_\xi\) gives the
corresponding virtual mass-addition rate per unit current mixture volume; no
additional mass variable is introduced.

Variations of the phase motion also change the volume ratio between the phase
reference and current configurations.  The compatible virtual Jacobian rate is
%
\begin{equation}
	\delta\dot J_\xi
	=
	J_\xi\nabla_{\mbf x}\cdot\delta\mbf v_\xi.
	\label{eq:vp_virtual_jacobian_rate}
\end{equation}
%
This follows by linearizing
\(\dot J_\xi=J_\xi\nabla_{\mbf x}\cdot\mbf v_\xi\) at the current state.

\subsection{Inertial virtual power}

The phase kinetic energy is first written on the phase reference configuration,
%
\begin{align}
	T_\xi
	&=
	\int_{\mc B_{\xi0}}
	\frac{1}{2}
	\sum_\alpha
	\left(
	\rho_{\xi0}^\alpha
	+
	\mc C_\xi^\alpha
	+
	\mc D_\xi^\alpha
	\right)
	\mbf v_\xi\cdot\mbf v_\xi\,\dVxi.
	\label{eq:vp_reference_kinetic_energy}
\end{align}
%
The first variation with respect to the phase velocity and accumulated
reference component masses is
%
\begin{align}
	\delta T_\xi
	={}&
	\int_{\mc B_{\xi0}}
	\sum_\alpha
	\left(
	\rho_{\xi0}^\alpha+\mc C_\xi^\alpha+\mc D_\xi^\alpha
	\right)
	\mbf v_\xi\cdot\delta\mbf v_\xi\,\dVxi
	\notag\\
	&+
	\int_{\mc B_{\xi0}}
	\frac{1}{2}\mbf v_\xi\cdot\mbf v_\xi
	\sum_\alpha
	\left(
	\delta\mc C_\xi^\alpha+\delta\mc D_\xi^\alpha
	\right)\dVxi.
	\label{eq:vp_reference_kinetic_differential}
\end{align}
%
The momentum conjugate to the phase velocity is therefore
%
\begin{equation}
	\sum_\alpha
	\left(
	\rho_{\xi0}^\alpha+\mc C_\xi^\alpha+\mc D_\xi^\alpha
	\right)\mbf v_\xi
	=
	J_\xi\rho_\xi\mbf v_\xi.
	\label{eq:vp_reference_phase_momentum}
\end{equation}
%
Taking its material time derivative gives
%
\begin{align}
	\frac{D_\xi}{Dt}
	\left(J_\xi\rho_\xi\mbf v_\xi\right)
	={}&
	J_\xi\rho_\xi\mbf a_\xi
	+
	\sum_\alpha
	\left(
	\dot{\mc C}_\xi^\alpha+\dot{\mc D}_\xi^\alpha
	\right)\mbf v_\xi.
	\label{eq:vp_reference_momentum_rate}
\end{align}
%
The inertial virtual power is the negative momentum-rate pairing together with
the kinetic-energy pairing conjugate to virtual mass addition.  Pushing
\eqref{eq:vp_reference_momentum_rate} forward with
\(\dv=J_\xi\dVxi\) gives
%
\begin{align}
	\delta\mc P_{\rm in}
	={}&
	\sum_\xi\int_{\mc B}
	\left[
		-\rho_\xi\mbf a_\xi
		-
		\sum_\alpha
		\frac{\dot{\mc C}_\xi^\alpha+\dot{\mc D}_\xi^\alpha}{J_\xi}
		\mbf v_\xi
	\right]\cdot\delta\mbf v_\xi\,\dv
	\notag\\
	&+
	\sum_\xi\sum_\alpha\int_{\mc B}
	\frac{1}{2}\mbf v_\xi\cdot\mbf v_\xi
	\frac{
		\delta\mc C_\xi^\alpha+\delta\mc D_\xi^\alpha
	}{J_\xi}\,\dv.
	\label{eq:vp_inertial_power}
\end{align}
%
Thus the open-system inertia and the kinetic contribution conjugate to component
mass addition follow from the reference kinetic energy; neither is postulated.

\subsection{Reversible internal virtual power}

Temperature is included as a constitutive state variable but is not varied in
the virtual-power principle.  The mechanical and compositional virtual powers
are evaluated at fixed temperature; the actual temperature rate and the heat
flux enter later through the energy balance and the entropy inequality.  At
fixed temperature, let the specific phase Helmholtz free energy have the
constitutive dependence
%
\begin{equation}
	\psi_\xi
	=
	\psi_\xi
	\left(
	\mbf F_\xi,
	\bar\rho_\xi,
	\left\{\eta_\xi^\alpha\right\},
	\theta
	\right).
	\label{eq:vp_phase_free_energy_state}
\end{equation}
%
Its virtual material rate is
%
\begin{align}
	\delta\dot\psi_\xi
	={}&
	\left.
	\frac{\partial\psi_\xi}{\partial\mbf F_\xi}
	\right\vert_{\bar\rho_\xi,\eta_\xi^\beta,\theta}
	:
	\left(
	\nabla_{\mbf x}\delta\mbf v_\xi\,\mbf F_\xi
	\right)
	\notag\\
	&+
	\left.
	\frac{\partial\psi_\xi}{\partial\bar\rho_\xi}
	\right\vert_{\mbf F_\xi,\eta_\xi^\beta,\theta}
	\delta\dot{\bar\rho}_\xi
	+
	\sum_\alpha
	\left.
	\frac{\partial\psi_\xi}{\partial\eta_\xi^\alpha}
	\right\vert_{\mbf F_\xi,\bar\rho_\xi,
	\eta_\xi^{\beta\ne\alpha},\theta}
	\delta\dot\eta_\xi^\alpha.
	\label{eq:vp_free_energy_virtual_rate}
\end{align}
%
Define the phase effective Cauchy stress by
%
\begin{equation}
	\bs\sigma_\xi'
	\equiv
	\rho_\xi
	\left.
	\frac{\partial\psi_\xi}{\partial\mbf F_\xi}
	\right\vert_{\bar\rho_\xi,\eta_\xi^\beta,\theta}
	\mbf F_\xi^T.
	\label{eq:vp_effective_stress_definition}
\end{equation}
%
The phase contribution to reversible internal virtual power is therefore
%
\begin{align}
	\delta\mc P_{\rm phase}
	={}&
	-\sum_\xi\int_{\mc B}
	\bs\sigma_\xi':\nabla_{\mbf x}\delta\mbf v_\xi\,\dv
	\notag\\
	&-
	\sum_\xi\int_{\mc B}
	\rho_\xi
	\left.
	\frac{\partial\psi_\xi}{\partial\bar\rho_\xi}
	\right\vert_{\mbf F_\xi,\eta_\xi^\beta,\theta}
	\delta\dot{\bar\rho}_\xi\,\dv
	\notag\\
	&-
	\sum_\xi\sum_\alpha\int_{\mc B}
	\rho_\xi
	\left.
	\frac{\partial\psi_\xi}{\partial\eta_\xi^\alpha}
	\right\vert_{\mbf F_\xi,\bar\rho_\xi,
	\eta_\xi^{\beta\ne\alpha},\theta}
	\delta\dot\eta_\xi^\alpha\,\dv
	\notag\\
	&-
	\sum_\xi\sum_\alpha\int_{\mc B}
	\psi_\xi
	\frac{
		\delta\mc C_\xi^\alpha+\delta\mc D_\xi^\alpha
	}{J_\xi}\,\dv.
	\label{eq:vp_phase_internal_power_before_parts}
\end{align}
%
Spatial integration by parts of the stress term gives
%
\begin{align}
	-\int_{\mc B}
	\bs\sigma_\xi':\nabla_{\mbf x}\delta\mbf v_\xi\,\dv
	={}&
	\int_{\mc B}
	\nabla_{\mbf x}\cdot\bs\sigma_\xi'
	\cdot\delta\mbf v_\xi\,\dv
	-
	\int_{\partial\mc B}
	\left(\bs\sigma_\xi'\mbf n\right)
	\cdot\delta\mbf v_\xi\,\mathrm da.
	\label{eq:vp_stress_integration_by_parts}
\end{align}
%
The boundary term supplies the effective-stress traction and is omitted from the
local field equations collected in \eqref{eq:vp_collected_principle}.

\subsection{Variation of internal surface energy}

For the interfacial energy, define the total fluid fraction and saturations by
%
\begin{equation}
	\phi
	=
	\sum_{f\in\mc F}\phi_f,
	\qquad
	S_f
	=
	\frac{\phi_f}{\phi},
	\qquad
	\sum_{f\in\mc F}S_f=1.
	\label{eq:vp_fluid_saturation_definitions}
\end{equation}
%
The virtual saturation rates are
%
\begin{equation}
	\delta\dot S_f
	=
	\frac{
	\delta\dot\phi_f-S_f\sum_{g\in\mc F}\delta\dot\phi_g
	}{\phi}.
	\label{eq:vp_virtual_saturation_rate}
\end{equation}
%
Let the interfacial energy density be
%
\begin{equation}
	\gamma
	=
	\gamma
	\left(
	\left\{S_f\right\}_{f\in\mc F},\theta
	\right).
	\label{eq:vp_interfacial_energy_density}
\end{equation}
%
Its derivative with respect to the primitive fluid volume fractions is
%
\begin{equation}
	\gamma_f
	=
	\gamma
	+
	\left.
	\frac{\partial\gamma}{\partial S_f}
	\right\vert_{S_{g\ne f}}
	-
	\sum_{g\in\mc F}S_g
	\left.
	\frac{\partial\gamma}{\partial S_g}
	\right\vert_{S_{h\ne g}}.
	\label{eq:vp_interfacial_phase_derivative}
\end{equation}
%
Using this simplex derivative and transporting each phase-attached volume
fraction by its own virtual velocity gives
%
\begin{align}
	\delta\mc P_\gamma
	={}&
	-\int_{\mc B}\sum_{f\in\mc F}
	\gamma_f\delta\dot\phi_f\,\dv
	+
	\int_{\mc B}\sum_{f\in\mc F}
	\gamma_f\nabla_{\mbf x}\phi_f
	\cdot\delta\mbf v_f\,\dv.
	\label{eq:vp_interfacial_internal_power}
\end{align}
%
The total reversible contribution is
\(\delta\mc P_{\rm rev}=\delta\mc P_{\rm phase}+\delta\mc P_\gamma\).

\subsection{External and dissipative virtual powers}

Let \(\mbf g\) denote gravitational acceleration, \(\mbf b_\xi\) the momentum
supply to phase \(\xi\) from interphase interactions, and \(L_\xi^\alpha\) the
generalized work potential conjugate to accumulated component mass.  Their
external virtual power is
%
\begin{equation}
	\delta\mc P_{\rm ext}
	=
	\sum_\xi\int_{\mc B}
	\left(
	\rho_\xi\mbf g+\mbf b_\xi
	\right)\cdot\delta\mbf v_\xi\,\dv
	-
	\sum_\xi\sum_\alpha\int_{\mc B}
	L_\xi^\alpha
	\frac{
		\delta\mc C_\xi^\alpha+\delta\mc D_\xi^\alpha
	}{J_\xi}\,\dv.
	\label{eq:vp_external_power}
\end{equation}
%
Let \(\mc A_{(m)}\) denote the generalized force conjugate to mechanism rate
\(\dot r_{(m)}\), and let \(\mbf M_{\xi,{\rm disp}}^\alpha\) and
\(\mbf M_{\xi,{\rm diff}}^\alpha\) denote the forces conjugate to the two
relative component fluxes.  The dissipative virtual power is
%
\begin{align}
	\delta\mc P_{\rm diss}
	={}&
	-\sum_m\int_{\mc B}
	\mc A_{(m)}\delta\dot r_{(m)}\,\dv
	\notag\\
	&-
	\sum_\xi\sum_\alpha\int_{\mc B}
	\mbf M_{\xi,{\rm disp}}^\alpha\cdot
	\delta\mbf j_{\xi,{\rm disp}}^\alpha\,\dv
	\notag\\
	&-
	\sum_\xi\sum_\alpha\int_{\mc B}
	\mbf M_{\xi,{\rm diff}}^\alpha\cdot
	\delta\mbf j_{\xi,{\rm diff}}^\alpha\,\dv.
	\label{eq:vp_dissipative_power}
\end{align}

\subsection{Virtual power of the volume-fraction constraint}

The transport identity needed for the volume-fraction and mass-fraction
constraints is derived once here.  Pulling the current volume fraction back to
the phase reference configuration gives the reference factor
%
\begin{equation*}
	\phi_\xi\left(\mbf x_\xi(\mbf X_\xi,t),t\right)J_\xi.
\end{equation*}
%
Taking its material time derivative and using
\(\dot J_\xi=J_\xi\nabla_{\mbf x}\cdot\mbf v_\xi\) gives
%
\begin{align*}
	\frac{D_\xi}{Dt}\left(\phi_\xi J_\xi\right)
	={}&
	J_\xi
	\left(
	\dot\phi_\xi
	+
	\phi_\xi\nabla_{\mbf x}\cdot\mbf v_\xi
	\right),
	\\
	\dot\phi_\xi
	={}&
	\left.\frac{\partial\phi_\xi}{\partial t}\right\vert_{\mbf x}
	+
	\nabla_{\mbf x}\phi_\xi\cdot\mbf v_\xi.
\end{align*}
%
Thus the fixed-spatial rate of the pointwise saturation constraint
\(\sum_\xi\phi_\xi=1\) is
%
\begin{equation*}
	0
	=
	\sum_\xi
	\left.\frac{\partial\phi_\xi}{\partial t}\right\vert_{\mbf x}
	=
	\sum_\xi
	\left(
	\dot\phi_\xi
	-
	\nabla_{\mbf x}\phi_\xi\cdot\mbf v_\xi
	\right).
\end{equation*}
%
At the virtual-rate level, vary the pulled-back factor
\(\delta\left(\phi_\xi J_\xi\right)\) and use
\(\delta\dot J_\xi=J_\xi\nabla_{\mbf x}\cdot\delta\mbf v_\xi\).  Substituting
this compatible virtual motion into the fixed-spatial rate identity gives
%
\begin{equation}
	\sum_\xi
	\left(
	\delta\dot\phi_\xi
	-
	\nabla_{\mbf x}\phi_\xi\cdot\delta\mbf v_\xi
	\right)
	=
	\sum_\xi
	\left.\delta\dot\phi_\xi\right\vert_{\mbf x}
	=0.
	\label{eq:vp_volume_fraction_virtual_constraint}
\end{equation}
%
The augmented virtual power is therefore
%
\begin{align}
	\delta\mc P_{C_1}
	={}&
	\int_{\mc B}\delta\lambda
	\sum_\xi
	\left(
	\dot\phi_\xi
	-
	\nabla_{\mbf x}\phi_\xi\cdot\mbf v_\xi
	\right)\dv
	\notag\\
	&+
	\sum_\xi\int_{\mc B}\lambda
	\left(
	\delta\dot\phi_\xi
	-
	\nabla_{\mbf x}\phi_\xi\cdot\delta\mbf v_\xi
	\right)\dv.
	\label{eq:vp_C1_power}
\end{align}
%
The multiplier variation preserves \(\sum_\xi\phi_\xi=1\) from an admissible
initial state.  The remaining terms contribute \(\lambda\) to the virtual
volume-fraction rate and \(-\lambda\nabla_{\mbf x}\phi_\xi\) to the virtual
phase velocity.

\subsection{Virtual power of the component material-measure constraint}

Taking the material time derivative of
\eqref{eq:vp_reference_component_measure}, using
\(\rho_\xi^\alpha=\phi_\xi\bar\rho_\xi\eta_\xi^\alpha\), and dividing by
\(J_\xi\phi_\xi\bar\rho_\xi\eta_\xi^\alpha\) gives the rate identity
%
\begin{align}
	0={}&
	\nabla_{\mbf x}\cdot\mbf v_\xi
	-
	\frac{
	\dot{\mc C}_\xi^\alpha+\dot{\mc D}_\xi^\alpha
	}{J_\xi\phi_\xi\bar\rho_\xi\eta_\xi^\alpha}
	+
	\frac{\dot\phi_\xi}{\phi_\xi}
	+
	\frac{\dot{\bar\rho}_\xi}{\bar\rho_\xi}
	+
	\frac{\dot\eta_\xi^\alpha}{\eta_\xi^\alpha}.
	\label{eq:vp_component_measure_rate_constraint}
\end{align}
%
Its virtual rate is obtained term by term as
%
\begin{align}
	0={}&
	\nabla_{\mbf x}\cdot\delta\mbf v_\xi
	-
	\frac{
		\delta\mc C_\xi^\alpha+\delta\mc D_\xi^\alpha
	}{J_\xi\phi_\xi\bar\rho_\xi\eta_\xi^\alpha}
	+
	\frac{\delta\dot\phi_\xi}{\phi_\xi}
	+
	\frac{\delta\dot{\bar\rho}_\xi}{\bar\rho_\xi}
	+
	\frac{\delta\dot\eta_\xi^\alpha}{\eta_\xi^\alpha}.
	\label{eq:vp_component_measure_virtual_rate}
\end{align}
%
Multiplication by \(\pi_\xi^\alpha\), summation, and integration give
%
\begin{align}
	\delta\mc P_{C_2}
	={}&
	\sum_\xi\sum_\alpha\int_{\mc B}
	\delta\pi_\xi^\alpha
	\left[
	\nabla_{\mbf x}\cdot\mbf v_\xi
	-
	\frac{
	\dot{\mc C}_\xi^\alpha+\dot{\mc D}_\xi^\alpha
	}{J_\xi\phi_\xi\bar\rho_\xi\eta_\xi^\alpha}
	+
	\frac{\dot\phi_\xi}{\phi_\xi}
	+
	\frac{\dot{\bar\rho}_\xi}{\bar\rho_\xi}
	+
	\frac{\dot\eta_\xi^\alpha}{\eta_\xi^\alpha}
	\right]\dv
	\notag\\
	&+
	\sum_\xi\sum_\alpha\int_{\mc B}
	\pi_\xi^\alpha
	\left[
	\nabla_{\mbf x}\cdot\delta\mbf v_\xi
	+
	\frac{\delta\dot\phi_\xi}{\phi_\xi}
	+
	\frac{\delta\dot{\bar\rho}_\xi}{\bar\rho_\xi}
	+
	\frac{\delta\dot\eta_\xi^\alpha}{\eta_\xi^\alpha}
	\right.\notag\\
	&\left.\hspace{42mm}
	-
	\frac{
		\delta\mc C_\xi^\alpha+\delta\mc D_\xi^\alpha
	}{J_\xi\phi_\xi\bar\rho_\xi\eta_\xi^\alpha}
	\right]\dv.
	\label{eq:vp_C2_power_before_parts}
\end{align}
%
Spatial integration by parts yields
%
\begin{align}
	\sum_\alpha\int_{\mc B}
	\pi_\xi^\alpha
	\nabla_{\mbf x}\cdot\delta\mbf v_\xi\,\dv
	={}&
	-\int_{\mc B}
	\sum_\alpha\nabla_{\mbf x}\pi_\xi^\alpha
	\cdot\delta\mbf v_\xi\,\dv
	+
	\text{b.t.}
	\label{eq:vp_C2_integration_by_parts}
\end{align}
%
This derives, rather than assumes, the storage-multiplier gradient in the phase
momentum equation.  The multiplier variation preserves
\eqref{eq:vp_reference_component_measure} from an admissible initial state.

\subsection{Virtual power of the mass-fraction constraint}

Applying the pull-back identity in
\eqref{eq:vp_volume_fraction_virtual_constraint} to each phase-attached mass
fraction gives the virtual rate of
\(\sum_\alpha\eta_\xi^\alpha=1\),
%
\begin{equation}
	\sum_\alpha
	\left(
	\delta\dot\eta_\xi^\alpha
	-
	\nabla_{\mbf x}\eta_\xi^\alpha\cdot\delta\mbf v_\xi
	\right)
	=0.
	\label{eq:vp_mass_fraction_virtual_constraint}
\end{equation}
%
Because \(\sum_\alpha\nabla_{\mbf x}\eta_\xi^\alpha=\mbf0\), its augmented
virtual power reduces to
%
\begin{align}
	\delta\mc P_{C_3}
	={}&
	\sum_\xi\int_{\mc B}\delta\Pi_\xi
	\sum_\alpha
	\left(
	\dot\eta_\xi^\alpha
	-
	\nabla_{\mbf x}\eta_\xi^\alpha\cdot\mbf v_\xi
	\right)\dv
	+
	\sum_\xi\sum_\alpha\int_{\mc B}
	\Pi_\xi\delta\dot\eta_\xi^\alpha\,\dv.
	\label{eq:vp_C3_power}
\end{align}
%
The multiplier variation preserves \(\sum_\alpha\eta_\xi^\alpha=1\) from an
admissible initial state.

\subsection{Virtual power of the component-rate constraint}

The complete current-volume component-rate constraint is
%
\begin{align}
	0={}&
	\frac{\dot{\mc C}_\xi^\alpha+\dot{\mc D}_\xi^\alpha}{J_\xi}
	-
	\sum_m\nu_{\xi (m)}^\alpha\dot r_{(m)}
	+
	\nabla_{\mbf x}\cdot
	\left(
	\mbf j_{\xi,{\rm disp}}^\alpha
	+
	\mbf j_{\xi,{\rm diff}}^\alpha
	\right).
	\label{eq:vp_complete_component_rate_constraint}
\end{align}
%
The displacement coupling cannot be omitted because the two accumulation rates
are referential.  For example, their multiplier pairing pulls back as
%
\begin{align}
	\int_{\mc B}
	\mu_\xi^\alpha
	\frac{\dot{\mc C}_\xi^\alpha+\dot{\mc D}_\xi^\alpha}{J_\xi}\,\dv
	={}&
	\int_{\mc B_{\xi0}}
	\mu_\xi^\alpha\left(\mbf x_\xi(\mbf X_\xi,t),t\right)
	\left(
	\dot{\mc C}_\xi^\alpha+\dot{\mc D}_\xi^\alpha
	\right)\dVxi.
	\label{eq:vp_component_rate_pullback}
\end{align}
%
Its variation contains the compatible accumulation-rate variation
%
\begin{align}
	\delta\left(
	\int_{\mc B_{\xi0}}
	\mu_\xi^\alpha
	\left(
	\dot{\mc C}_\xi^\alpha+\dot{\mc D}_\xi^\alpha
	\right)\dVxi
	\right)
	={}&
	\int_{\mc B}
	\mu_\xi^\alpha
	\frac{
	\delta\dot{\mc C}_\xi^\alpha+\delta\dot{\mc D}_\xi^\alpha
	}{J_\xi}\,\dv
	\notag\\
	&+
	\int_{\mc B}
	\frac{\dot{\mc C}_\xi^\alpha+\dot{\mc D}_\xi^\alpha}{J_\xi}
	\nabla_{\mbf x}\mu_\xi^\alpha
	\cdot\delta\mbf v_\xi\,\dv.
	\label{eq:vp_component_rate_pullback_variation}
\end{align}
%
The first term in \eqref{eq:vp_component_rate_pullback_variation} must be
integrated by parts in time.  Because the accumulated-mass variations vanish at
\(t_1\) and \(t_2\),
%
\begin{align*}
	&\int_{t_1}^{t_2}\int_{\mc B}
	\mu_\xi^\alpha
	\frac{
	\delta\dot{\mc C}_\xi^\alpha+
	\delta\dot{\mc D}_\xi^\alpha
	}{J_\xi}\,\dv\dt
	\\
	&\qquad={}
	-\int_{t_1}^{t_2}\int_{\mc B}
	\frac{D_\xi\mu_\xi^\alpha}{Dt}
	\frac{
	\delta\mc C_\xi^\alpha+
	\delta\mc D_\xi^\alpha
	}{J_\xi}\,\dv\dt.
\end{align*}
%
Thus the accumulated masses produce the material derivative of the component-rate
multiplier.  The second term in
\eqref{eq:vp_component_rate_pullback_variation} is the source of the
multiplier-gradient contribution to momentum; it is not inserted as an
assumed body force.

The flux contribution is integrated by parts before taking its variation,
%
\begin{align}
	\int_{\mc B}\mu_\xi^\alpha
	\nabla_{\mbf x}\cdot\mbf j_{\xi,k}^\alpha\,\dv
	={}&
	-\int_{\mc B}
	\nabla_{\mbf x}\mu_\xi^\alpha\cdot
	\mbf j_{\xi,k}^\alpha\,\dv
	+
	\text{b.t.},
	\qquad
	k\in\left\{{\rm disp},{\rm diff}\right\}.
	\label{eq:vp_flux_constraint_parts}
\end{align}
%
Consequently its process-rate variation is
%
\begin{equation}
	-\int_{\mc B}
	\nabla_{\mbf x}\mu_\xi^\alpha\cdot
	\delta\mbf j_{\xi,k}^\alpha\,\dv.
	\label{eq:vp_flux_constraint_variation}
\end{equation}
%
After the temporal integration by parts, collecting the multiplier,
accumulation, motion, mechanism-rate, and flux variations gives
%
\begin{align}
	\delta\mc P_{C_4}
	={}&
	\sum_\xi\sum_\alpha\int_{\mc B}
	\delta\mu_\xi^\alpha
	\left[
	\frac{\dot{\mc C}_\xi^\alpha+\dot{\mc D}_\xi^\alpha}{J_\xi}
	-
	\sum_m\nu_{\xi (m)}^\alpha\dot r_{(m)}
	\right.\notag\\
	&\left.\hspace{36mm}
	+
	\nabla_{\mbf x}\cdot
	\left(
	\mbf j_{\xi,{\rm disp}}^\alpha
	+
	\mbf j_{\xi,{\rm diff}}^\alpha
	\right)
	\right]\dv
	\notag\\
	&-
	\sum_\xi\sum_\alpha\int_{\mc B}
	\frac{D_\xi\mu_\xi^\alpha}{Dt}
	\frac{
	\delta\mc C_\xi^\alpha+\delta\mc D_\xi^\alpha
	}{J_\xi}\,\dv
	\notag\\
	&+
	\sum_\xi\sum_\alpha\int_{\mc B}
	\frac{\dot{\mc C}_\xi^\alpha+\dot{\mc D}_\xi^\alpha}{J_\xi}
	\nabla_{\mbf x}\mu_\xi^\alpha
	\cdot\delta\mbf v_\xi\,\dv
	\notag\\
	&-
	\sum_m\int_{\mc B}
	\left(
	\sum_\xi\sum_\alpha
	\mu_\xi^\alpha\nu_{\xi (m)}^\alpha
	\right)\delta\dot r_{(m)}\,\dv
	\notag\\
	&-
	\sum_\xi\sum_\alpha\int_{\mc B}
	\nabla_{\mbf x}\mu_\xi^\alpha\cdot
	\left(
	\delta\mbf j_{\xi,{\rm disp}}^\alpha
	+
	\delta\mbf j_{\xi,{\rm diff}}^\alpha
	\right)\dv
	+
	\text{b.t.}
	\label{eq:vp_C4_power}
\end{align}

\subsection{Electric enthalpy and the component-rate multiplier}

The charge bookkeeping in \sref{conservation_of_charge} supplies the source
term for the electric field.  In the electroquasistatic specialization the
corresponding reversible field variable is the electric potential
\(\varphi\).  Introduce the electric field \(\mbf E\), electric displacement
\(\mbf d\), and, when needed, a polarization \(\mbf P\),
%
\begin{equation}
	\mbf E=-\nabla_{\mbf x}\varphi,
	\qquad
	\mbf d
	=
	-\left.
	\frac{\partial\omega_{\rm elec}}{\partial\mbf E}
	\right\vert_{\mbf P,\mc X_\xi,\varphi}.
	\label{eq:electric_field_definitions}
\end{equation}
%
Here \(\omega_{\rm elec}\) is an electric enthalpy per unit current mixture
volume.  The notation \(\mc X_\xi\) denotes the mechanical, compositional, and
thermal state variables already used by the phase Helmholtz energies.  A linear
dielectric without an independent polarization is recovered by choosing
\(\omega_{\rm elec}=-\frac{1}{2}\mbf E\cdot\bs\epsilon\mbf E\), which gives
\(\mbf d=\bs\epsilon\mbf E\).  In reference-form solid calculations this
current-volume term is pulled back as \(J\omega_{\rm elec}\), with the same
Jacobian used for the current-volume virtual power.

The reversible electrical contribution is
%
\begin{equation}
	\mc I_{\rm elec}[\varphi,\mbf P]
	=
	\int_{\mc B}
	\left[
	\omega_{\rm elec}
	\left(
	\mbf E,\mbf P,\nabla_{\mbf x}\mbf P,\mc X_\xi
	\right)
	+
	\varrho\varphi
	\right]\dv ,
	\qquad
	\mbf E=-\nabla_{\mbf x}\varphi .
	\label{eq:electric_enthalpy_functional}
\end{equation}
%
To obtain the potential equation in the same rate form as the rest of the
virtual-power principle, differentiate \eqref{eq:electric_enthalpy_functional}
in time and isolate the terms proportional to \(\dot\varphi\).  Terms involving
the rates of \(\varrho\), \(\mbf P\), the other state variables, and the motion
of the current domain belong to their own conjugate rate coefficients and do not
change the coefficient of \(\dot\varphi\).  Since
\(\dot{\mbf E}=-\nabla_{\mbf x}\dot\varphi\) and
\(\mbf d=-\partial\omega_{\rm elec}/\partial\mbf E\), the potential-rate
contribution is
%
\begin{align*}
	{}&
	\int_{\mc B}
	\left[
	\left.
	\frac{\partial\omega_{\rm elec}}{\partial\mbf E}
	\right\vert_{\mbf P,\mc X_\xi,\varphi}
	\cdot\dot{\mbf E}
	+
	\varrho\dot\varphi
	\right]\dv
	\\
	={}&
	\int_{\mc B}
	\left[
	\mbf d\cdot\nabla_{\mbf x}\dot\varphi
	+
	\varrho\dot\varphi
	\right]\dv
	\\
	={}&
	\int_{\mc B}
	\left(
	\varrho-\nabla_{\mbf x}\cdot\mbf d
	\right)\dot\varphi\,\dv
	+
	\int_{\partial\mc B}
	\left(
	\mbf d\cdot\mbf n
	\right)\dot\varphi\,da .
\end{align*}
%
The corresponding virtual-power term is obtained by pairing this same interior
coefficient with the admissible virtual potential rate \(\delta\dot\varphi\):
%
\begin{align}
	\delta\mc P_{\rm elec}^{\varphi}
	={}&
	\int_{\mc B}
	\left(
	\varrho-\nabla_{\mbf x}\cdot\mbf d
	\right)\delta\dot\varphi\,\dv
	+
	\text{b.t.},
	\label{eq:vp_electric_potential_variation}
\end{align}
%
so the coefficient of the arbitrary interior virtual rate
\(\delta\dot\varphi\) gives Gauss' law,
%
\begin{equation}
	\nabla_{\mbf x}\cdot\mbf d=\varrho .
	\label{eq:gauss_law}
\end{equation}
%
Thus Gauss' law is the Euler--Lagrange equation associated with the reversible
electric enthalpy.  It is distinct from the dielectric closure, which is the
constitutive choice of \(\omega_{\rm elec}\), or equivalently of \(\mbf d\), as
a function of electric, mechanical, compositional, and thermal state.  The
natural electrical boundary condition is prescribed normal displacement
\(\mbf d\cdot\mbf n\), while the essential boundary condition is prescribed
\(\varphi\).

The same rate collection gives the polarization contribution
%
\begin{align}
	\delta\mc P_{\rm elec}^{\mbf P}
	={}&
	\int_{\mc B}
	\left[
	-\left.
	\frac{\partial\omega_{\rm elec}}{\partial\mbf P}
	\right\vert_{\mbf E,\nabla_{\mbf x}\mbf P,\mc X_\xi}
	+
	\nabla_{\mbf x}\cdot
	\left.
	\frac{\partial\omega_{\rm elec}}
	{\partial\nabla_{\mbf x}\mbf P}
	\right\vert_{\mbf E,\mbf P,\mc X_\xi}
	\right]\cdot\delta\dot{\mbf P}\,\dv
	+
	\text{b.t.}
	\label{eq:vp_polarization_variation}
\end{align}
%
If polarization is an independent reversible field, stationarity with respect to
\(\mbf P\) gives
%
\begin{equation}
	\left.
	\frac{\partial\omega_{\rm elec}}{\partial\mbf P}
	\right\vert_{\mbf E,\nabla_{\mbf x}\mbf P,\mc X_\xi}
	-
	\nabla_{\mbf x}\cdot
	\left.
	\frac{\partial\omega_{\rm elec}}
	{\partial\nabla_{\mbf x}\mbf P}
	\right\vert_{\mbf E,\mbf P,\mc X_\xi}
	=
	\mbf0 .
	\label{eq:polarization_stationarity}
\end{equation}
%
If polarization switching, dielectric relaxation, or interfacial polarization is
dissipative, this stationarity condition is not imposed; the same coefficient is
retained as the force conjugate to \(\dot{\mbf P}\) in the residual entropy
inequality.

The electroquasistatic field equation \eqref{eq:gauss_law} and the local charge
balance \eqref{eq:mixture_charge_balance} are compatible because
%
\begin{equation}
	\nabla_{\mbf x}\cdot
	\left(
	\dot{\mbf d}
	+
	\mbf i
	\right)
	=0 .
	\label{eq:charge_gauss_compatibility}
\end{equation}
%
The solenoidal part of \(\dot{\mbf d}+\mbf i\) is not resolved in this
electroquasistatic reduction.  A full electromagnetic specialization would add a
vector potential, magnetic induction, Faraday's law, Ampere--Maxwell balance,
and the corresponding electromagnetic momentum and energy fluxes.  The scalar
potential theory used here is the first specialization for ion exchange,
double-layer, electro-osmotic, and electrochemical porous-flow models in which
magnetic induction is not a primary state variable.

The same electrical functional contributes to the coefficient of each charged
component material measure because the charge carried by a component is
\(z_\xi^\alpha\) times that component's mass.  Written against the
charge-weighted rate identity \eqref{eq:component_charge_rate_identity}, this
contribution has the form
%
\begin{align}
	0={}&
	\sum_\xi\sum_\alpha
	\varphi
	\left[
	\frac{
	z_\xi^\alpha\dot{\mc C}_\xi^\alpha
	+
	z_\xi^\alpha\dot{\mc D}_\xi^\alpha
	}{J_\xi}
	-
	z_\xi^\alpha
	\sum_m\nu_{\xi (m)}^\alpha\dot r_{(m)}
	+
	z_\xi^\alpha
	\nabla_{\mbf x}\cdot
	\left(
	\mbf j_{\xi,{\rm disp}}^\alpha
	+
	\mbf j_{\xi,{\rm diff}}^\alpha
	\right)
	\right].
	\label{eq:vp_charge_weighted_component_rate_identity}
\end{align}
%
This is the component-rate constraint viewed through the charge weights
\(z_\xi^\alpha\).  Integrating the flux term by parts gives the local
electrical process power
%
\begin{align}
	{}&
	-\sum_m
	\left(
	\sum_\xi\sum_\alpha
	z_\xi^\alpha\varphi\nu_{\xi (m)}^\alpha
	\right)\delta\dot r_{(m)}
	\notag\\
	&-
	\sum_\xi\sum_\alpha
	z_\xi^\alpha\nabla_{\mbf x}\varphi\cdot
	\left(
	\delta\mbf j_{\xi,{\rm disp}}^\alpha
	+
	\delta\mbf j_{\xi,{\rm diff}}^\alpha
	\right)
	+
	\text{b.t.}
	\label{eq:vp_electrical_process_power}
\end{align}
%
Combining \eqref{eq:vp_electrical_process_power} with
\eqref{eq:vp_C4_power} shows that electrical work enters the same
component-rate multiplier \(\mu_\xi^\alpha\).  If a nonelectrical chemical part
is separated, then
\(\mu_\xi^\alpha=\mu_{\xi,{\rm chem}}^\alpha+z_\xi^\alpha\varphi\).  The
process-rate coefficients retain their original form,
%
\begin{align}
	\mc A_{(m)}
	&=
	-\sum_\xi\sum_\alpha
	\mu_\xi^\alpha\nu_{\xi (m)}^\alpha,
	\label{eq:vp_electrochemical_affinity_balance}
	\\
	\mbf M_{\xi,k}^\alpha
	&=
	-\nabla_{\mbf x}\mu_\xi^\alpha,
	\qquad
	k\in\left\{{\rm disp},{\rm diff}\right\}.
	\label{eq:vp_electrochemical_relative_flux_force_balance}
\end{align}
%
For a charge-conserving local mechanism satisfying
\eqref{eq:mechanism_charge_conservation},
the explicit \(\varphi\)-part cancels from the affinity if the decomposition
above is used.  The electric field then affects ion-exchange kinetics indirectly
through charged transport and through the state dependence of the chemical
potentials.

\subsection{Collected virtual-power principle}

Substitution of
\eqref{eq:vp_inertial_power},
\eqref{eq:vp_phase_internal_power_before_parts},
	\eqref{eq:vp_stress_integration_by_parts},
	\eqref{eq:vp_interfacial_internal_power},
	\eqref{eq:vp_electric_potential_variation},
	\eqref{eq:vp_polarization_variation},
	\eqref{eq:vp_electrical_process_power},
	\eqref{eq:vp_external_power},
	\eqref{eq:vp_dissipative_power}, and
\eqref{eq:vp_C1_power}--\eqref{eq:vp_C4_power} into
\eqref{eq:vp_principle} gives, after collecting independent virtual rates,
%
\begin{align}
	0={}&
	\sum_\xi\int_{\mc B}
	\left[
	-\rho_\xi\mbf a_\xi
	+
	\nabla_{\mbf x}\cdot\bs\sigma_\xi'
	+
	\rho_\xi\mbf g
	+
		\mbf b_\xi
		-
		\lambda\nabla_{\mbf x}\phi_\xi
		-
		\sum_\alpha\nabla_{\mbf x}\pi_\xi^\alpha
		+
		\sum_\alpha
		\dot m_\xi^\alpha
	\left(
	\nabla_{\mbf x}\mu_\xi^\alpha-\mbf v_\xi
	\right)
	\right]\cdot\delta\mbf v_\xi\,\dv
	\notag\\
	&+
	\sum_{f\in\mc F}\int_{\mc B}
	\gamma_f\nabla_{\mbf x}\phi_f\cdot\delta\mbf v_f\,\dv
	\notag\\
	&+
	\sum_\xi\int_{\mc B}
	\left[
	\lambda
	+
	\sum_\alpha\frac{\pi_\xi^\alpha}{\phi_\xi}
	\right]\delta\dot\phi_\xi\,\dv
	\notag\\
	&-
	\sum_{f\in\mc F}\int_{\mc B}
	\gamma_f\delta\dot\phi_f\,\dv
	\notag\\
	&+
	\sum_\xi\int_{\mc B}
	\left[
	-\rho_\xi
	\left.
	\frac{\partial\psi_\xi}{\partial\bar\rho_\xi}
	\right\vert_{\mbf F_\xi,\eta_\xi^\beta,\theta}
	+
	\sum_\alpha\frac{\pi_\xi^\alpha}{\bar\rho_\xi}
	\right]\delta\dot{\bar\rho}_\xi\,\dv
	\notag\\
	&+
	\sum_\xi\sum_\alpha\int_{\mc B}
	\left[
	\frac{\pi_\xi^\alpha}{\eta_\xi^\alpha}
	+
	\Pi_\xi
	-
	\rho_\xi
	\left.
	\frac{\partial\psi_\xi}{\partial\eta_\xi^\alpha}
	\right\vert_{\mbf F_\xi,\bar\rho_\xi,
	\eta_\xi^{\beta\ne\alpha},\theta}
	\right]\delta\dot\eta_\xi^\alpha\,\dv
	\notag\\
	&+
	\sum_\xi\sum_\alpha\int_{\mc B}
	\left[
	\frac{1}{2}\mbf v_\xi\cdot\mbf v_\xi
	-
	\psi_\xi
	-
	L_\xi^\alpha
		-
		\frac{D_\xi\mu_\xi^\alpha}{Dt}
		-
	\frac{\pi_\xi^\alpha}
	{\phi_\xi\bar\rho_\xi\eta_\xi^\alpha}
	\right]
	\frac{\delta\mc C_\xi^\alpha+\delta\mc D_\xi^\alpha}{J_\xi}
	\,\dv
	\notag\\
	&-
	\sum_m\int_{\mc B}
	\left[
	\mc A_{(m)}
	+
	\sum_\xi\sum_\alpha
	\mu_\xi^\alpha\nu_{\xi (m)}^\alpha
	\right]\delta\dot r_{(m)}\,\dv
	\notag\\
	&-
	\sum_\xi\sum_\alpha\int_{\mc B}
	\left(
	\mbf M_{\xi,{\rm disp}}^\alpha
	+
	\nabla_{\mbf x}\mu_\xi^\alpha
	\right)\cdot
	\delta\mbf j_{\xi,{\rm disp}}^\alpha\,\dv
	\notag\\
	&-
	\sum_\xi\sum_\alpha\int_{\mc B}
	\left(
	\mbf M_{\xi,{\rm diff}}^\alpha
	+
	\nabla_{\mbf x}\mu_\xi^\alpha
	\right)\cdot
	\delta\mbf j_{\xi,{\rm diff}}^\alpha\,\dv
	\notag\\
	&+
	\int_{\mc B}
	\left(
	\varrho-\nabla_{\mbf x}\cdot\mbf d
	\right)\delta\dot\varphi\,\dv
	\notag\\
	&+
	\int_{\mc B}
	\left[
	-\left.
	\frac{\partial\omega_{\rm elec}}{\partial\mbf P}
	\right\vert_{\mbf E,\nabla_{\mbf x}\mbf P,\mc X_\xi}
	+
	\nabla_{\mbf x}\cdot
	\left.
	\frac{\partial\omega_{\rm elec}}
	{\partial\nabla_{\mbf x}\mbf P}
	\right\vert_{\mbf E,\mbf P,\mc X_\xi}
	\right]\cdot\delta\dot{\mbf P}\,\dv
	\notag\\
	&+
	\text{multiplier variations}
	+
	\text{boundary powers}.
		\label{eq:vp_collected_principle}
\end{align}
%
For compactness in the balance equations below, define the total current-volume
component mass-addition rate by
%
\begin{equation*}
	\dot m_\xi^\alpha
	\equiv
	\frac{\dot{\mc C}_\xi^\alpha}{J_\xi}
	+
	\frac{\dot{\mc D}_\xi^\alpha}{J_\xi}
	=
	\dot c_\xi^\alpha
	+
	\frac{\dot{\mc D}_\xi^\alpha}{J_\xi}.
\end{equation*}
%
This notation combines conversion and relative-transport accumulation only
where they multiply the same mechanical or thermodynamic factor; the two rates
remain separate in their balance and dissipation equations.

The multiplier variations return the four constraints, while the coefficient of
the virtual potential rate returns Gauss' law.  If polarization is treated as a
reversible field, the coefficient of \(\delta\dot{\mbf P}\) gives
\eqref{eq:polarization_stationarity}; if it is treated as a dissipative internal
variable, the same coefficient is retained for the residual entropy inequality.
Because the remaining virtual rates and accumulated-state variations are
arbitrary, their coefficients vanish independently.
In particular, the phase-momentum coefficient contains
%
\begin{equation*}
	\sum_\alpha
		\dot m_\xi^\alpha
	\left(
	\nabla_{\mbf x}\mu_\xi^\alpha-\mbf v_\xi
	\right),
\end{equation*}
%
whose multiplier-gradient part was derived explicitly in
\eqref{eq:vp_component_rate_pullback}--
\eqref{eq:vp_component_rate_pullback_variation}.  The reaction and relative-flux
coefficients give
%
\begin{align}
	\mc A_{(m)}
	&=
	-\sum_\xi\sum_\alpha
	\mu_\xi^\alpha\nu_{\xi (m)}^\alpha,
	\label{eq:vp_affinity_balance}
	\\
	\mbf M_{\xi,k}^\alpha
	&=
	-\nabla_{\mbf x}\mu_\xi^\alpha,
	\qquad
	k\in\left\{{\rm disp},{\rm diff}\right\}.
	\label{eq:vp_relative_flux_force_balance}
\end{align}
%
With electrical effects included, the same equations are read with
\(\mu_\xi^\alpha\) as the electrochemical potential; this is the content of
\eqref{eq:vp_electrochemical_affinity_balance} and
\eqref{eq:vp_electrochemical_relative_flux_force_balance}.  These are
generalized-force balances rather than constitutive laws.  Admissible
constitutive relations between the generalized forces and their process rates
must make the corresponding reaction and relative-transport powers
nonnegative.  A quadratic dissipation potential then gives the linear Onsager
closures, while nonlinear convex dissipation potentials provide the analogous
farther-from-equilibrium rate laws~\cite{doi2011onsager,gurtin2010continua}.

\subsection{Field equations from virtual power}

The coefficient of each independent virtual rate in
\eqref{eq:vp_collected_principle} must vanish.  For a phase without an explicit
capillary contribution, the coefficient of \(\delta\dot\phi_\xi\) gives
%
\begin{equation}
	\phi_\xi\lambda
	=
	-\sum_\alpha\pi_\xi^\alpha.
	\label{eq:vp_el_lambda}
\end{equation}
%
The corresponding fluid-phase relation retains the interfacial contribution and
is obtained in \eqref{eq:vp_fluid_volume_fraction_relation}.

Using \eqref{eq:vp_el_lambda} in the coefficient of
\(\delta\mbf v_\xi\) gives the phase momentum equation
%
\begin{equation}
	\rho_\xi\mbf a_\xi
	=
	\nabla_{\mbf x}\cdot\bs\sigma_\xi'
	+
	\rho_\xi\mbf g
	+
	\mbf b_\xi
	+
	\phi_\xi\nabla_{\mbf x}\lambda
	+
	\sum_\alpha\dot m_\xi^\alpha
	\left(
	\nabla_{\mbf x}\mu_\xi^\alpha-\mbf v_\xi
	\right).
	\label{eq:vp_momentum_uniform}
\end{equation}
%
Thus both conversion and relative-transport accumulation contribute through the
same chemical-potential gradient and are combined in \(\dot m_\xi^\alpha\),
while their rates remain distinct in the component balance.

The coefficient of \(\delta\dot{\bar\rho}_\xi\) gives
%
\begin{equation}
	\sum_\alpha\frac{\pi_\xi^\alpha}{\phi_\xi}
	=
	\bar\rho_\xi^2
	\left.
	\frac{\partial\psi_\xi}{\partial\bar\rho_\xi}
	\right\vert_{\mbf F_\xi,\eta_\xi^\beta,\theta}.
	\label{eq:vp_phase_pressure_pi}
\end{equation}
%
Define the thermodynamic phase pressure by
%
\begin{equation}
	\boxed{
	\bar p_\xi
	\equiv
	\bar\rho_\xi^2
	\left.
	\frac{\partial\psi_\xi}{\partial\bar\rho_\xi}
	\right\vert_{\mbf F_\xi,\eta_\xi^\beta,\theta}
	}.
	\label{eq:vp_phase_pressure_definition}
\end{equation}
%
For a fluid phase, the coefficient of \(\delta\dot\phi_f\) in
\eqref{eq:vp_collected_principle} includes the primitive interfacial potential,
%
\begin{equation}
	\lambda
	+
	\sum_\alpha\frac{\pi_f^\alpha}{\phi_f}
	-
	\gamma_f
	=0.
	\label{eq:vp_fluid_volume_fraction_relation}
\end{equation}
%
Using \eqref{eq:vp_phase_pressure_definition} gives
%
\begin{equation}
	\lambda+\bar p_f-\gamma_f=0,
	\qquad
	\bar p_f=\gamma_f-\lambda.
	\label{eq:vp_fluid_pressure_primitive_relation}
\end{equation}
%
Consequently,
%
\begin{equation}
	\bar p_f-\bar p_g
	=
	\gamma_f-\gamma_g.
	\label{eq:vp_fluid_pressure_difference}
\end{equation}
%
Let \(\mc I\) be the set of oriented wetting/non-wetting fluid-phase pairs.
For \((w,n)\in\mc I\), define
%
\begin{equation}
	p_{c,nw}
	\equiv
	\gamma_n-\gamma_w.
	\label{eq:vp_capillary_potential_difference}
\end{equation}
%
Then
%
\begin{empheq}[box=\fbox]{equation}
	p_{c,nw}
	=
	\bar p_n-\bar p_w,
	\qquad
	(w,n)\in\mc I.
	\label{eq:vp_capillary_pressure_recovery}
\end{empheq}
%
For two fluid phases and an interfacial energy parameterized by the wetting
saturation, this reduces to
\(p_{c,nw}=-\partial\gamma/\partial S_w\).

The multiplier \(\lambda\) also determines the equivalent pore pressure.  Define
%
\begin{equation}
	\bar p_E
	\equiv
	-\lambda.
	\label{eq:vp_equivalent_pore_pressure_definition}
\end{equation}
%
Multiplying \(\bar p_f=\gamma_f+\bar p_E\) by \(S_f\), summing over the fluid
phases, and using \(\sum_{f\in\mc F}S_f=1\) gives
%
\begin{equation}
	\sum_{f\in\mc F}S_f\bar p_f
	=
	\bar p_E
	+
	\sum_{f\in\mc F}S_f\gamma_f.
	\label{eq:vp_saturation_weighted_pressure_sum}
\end{equation}
%
The definition \eqref{eq:vp_interfacial_phase_derivative} implies
%
\begin{equation}
	\sum_{f\in\mc F}S_f\gamma_f
	=
	\gamma.
	\label{eq:vp_saturation_weighted_interfacial_sum}
\end{equation}
%
Therefore,
%
\begin{equation}
	\boxed{
	\bar p_E
	=
	\sum_{f\in\mc F}S_f\bar p_f-\gamma
	}.
	\label{eq:vp_equivalent_pore_pressure_relation}
\end{equation}

For a phase with \(N\) tracked components, the coefficients of the independent
virtual composition rates give
%
\begin{equation}
	0
	=
	\frac{\pi_\xi^\alpha}{\eta_\xi^\alpha}
	+
	\Pi_\xi
	-
	\rho_\xi
	\left.
	\frac{\partial\psi_\xi}{\partial\eta_\xi^\alpha}
	\right\vert_{\mbf F_\xi,\bar\rho_\xi,
	\eta_\xi^{\beta\ne\alpha},\theta},
	\qquad
	\alpha=1,\ldots,N.
	\label{eq:vp_composition_stationarity}
\end{equation}
%
The multiplier variation preserves
%
\begin{equation}
	\sum_{\alpha=1}^N\eta_\xi^\alpha=1.
	\label{eq:vp_composition_constraint}
\end{equation}
%
Equation \eqref{eq:vp_composition_stationarity} can be written componentwise as
%
\begin{equation}
	\rho_\xi
	\left.
	\frac{\partial\psi_\xi}{\partial\eta_\xi^\alpha}
	\right\vert_{\mbf F_\xi,\bar\rho_\xi,
	\eta_\xi^{\beta\ne\alpha},\theta}
	-
	\frac{\pi_\xi^\alpha}{\eta_\xi^\alpha}
	=
	\Pi_\xi.
	\label{eq:vp_composition_multiplier_componentwise}
\end{equation}
%
Choosing component \(N\) as the reference and comparing it with
\(\beta=1,\ldots,N-1\) eliminates \(\Pi_\xi\),
%
\begin{align}
	\rho_\xi
	\left(
	\left.
	\frac{\partial\psi_\xi}{\partial\eta_\xi^\beta}
	\right\vert_{\mbf F_\xi,\bar\rho_\xi,
	\eta_\xi^{\kappa\ne\beta},\theta}
	-
	\left.
	\frac{\partial\psi_\xi}{\partial\eta_\xi^N}
	\right\vert_{\mbf F_\xi,\bar\rho_\xi,
	\eta_\xi^{\kappa\ne N},\theta}
	\right)
	={}&
	\frac{\pi_\xi^\beta}{\eta_\xi^\beta}
	-
	\frac{\pi_\xi^N}{\eta_\xi^N},
	\qquad
	\beta=1,\ldots,N-1.
	\label{eq:vp_composition_difference_system}
\end{align}
%
Here \(\kappa\) denotes the components held fixed in the indicated partial
derivative and is not summed.

The common coefficient of
\(\delta\mc C_\xi^\alpha/J_\xi\) and
\(\delta\mc D_\xi^\alpha/J_\xi\) gives the component-rate-multiplier evolution
equation
%
\begin{equation}
	\frac{D_\xi\mu_\xi^\alpha}{Dt}
	=
	\frac{1}{2}\mbf v_\xi\cdot\mbf v_\xi
	-
	\psi_\xi
	-
	L_\xi^\alpha
	-
	\frac{\pi_\xi^\alpha}
	{\phi_\xi\bar\rho_\xi\eta_\xi^\alpha}.
	\label{eq:vp_component_potential_evolution}
\end{equation}
%
Both accumulation mechanisms give the same equation because they enter the
reference mass measure through
\(\mc C_\xi^\alpha+\mc D_\xi^\alpha\), have the same generalized work
potential, and are governed by the same multiplier.  This is the
phase/component analogue of the Bedford--Drumheller multiplier evolution
equation~\cite[eq.~(48)]{drumheller2000}.  The process-rate
coefficients independently give the affinity and relative-flux force balances
\eqref{eq:vp_affinity_balance} and
\eqref{eq:vp_relative_flux_force_balance}.  These balances do not prescribe
kinetic or transport laws; admissible constitutive laws must satisfy the entropy
inequality.

\subsection{Fluid-phase momentum equations}
\label{sec:vp_fluid_phase_momentum}

For any fluid phase \(f\in\mc F\), retaining the scalar-transport capillary
contribution in \eqref{eq:vp_momentum_uniform} gives
%
\begin{align}
	\rho_f\mbf a_f
	={}&
	\nabla_{\mbf x}\cdot\bs\sigma_f'
	+
	\rho_f\mbf g
	+
	\mbf b_f
	-
	\lambda\nabla_{\mbf x}\phi_f
	-
	\sum_\alpha\nabla_{\mbf x}\pi_f^\alpha
	\notag\\
		&+
		\gamma_f\nabla_{\mbf x}\phi_f
		+
		\sum_\alpha\dot m_f^\alpha
		\left(
	\nabla_{\mbf x}\mu_f^\alpha-\mbf v_f
	\right).
	\label{eq:vp_fluid_momentum_intermediate}
\end{align}
%
Using
\(\sum_\alpha\pi_f^\alpha=\phi_f\bar p_f\) from
\eqref{eq:vp_phase_pressure_pi} and
\(\lambda+\bar p_f-\gamma_f=0\) from
\eqref{eq:vp_fluid_pressure_primitive_relation}, the multiplier and interfacial
gradient terms reduce as
%
\begin{align}
	-&\lambda\nabla_{\mbf x}\phi_f
	-
	\nabla_{\mbf x}\left(\phi_f\bar p_f\right)
	+
	\gamma_f\nabla_{\mbf x}\phi_f
	\notag\\
	&=
	-\phi_f\nabla_{\mbf x}\bar p_f.
	\label{eq:vp_fluid_pressure_gradient_collapse}
\end{align}
%
For a fluid free energy with no explicit dependence on \(\mbf F_f\),
\eqref{eq:vp_effective_stress_definition} gives \(\bs\sigma_f'=\mbf0\).  The
fluid-phase momentum equation is therefore
%
\begin{empheq}[box=\fbox]{equation}
	\rho_f\mbf a_f
	=
	\rho_f\mbf g
	+
	\mbf b_f
	-
	\phi_f\nabla_{\mbf x}\bar p_f
	+
	\sum_\alpha\dot m_f^\alpha
	\left(
	\nabla_{\mbf x}\mu_f^\alpha-\mbf v_f
	\right),
	\qquad f\in\mc F.
	\label{eq:vp_fluid_momentum_simplified}
\end{empheq}
